\section{SmallWood Programming Model}
\label{sec:smallwood-model}

This section specifies the SmallWood statements used by our ARC
instantiation. The abstract credential message is always
\[
  \vecm=\vecmu\Vert \veck,
\]
where \(\vecmu\) denotes the hidden attributes and \(\veck\) denotes the
hidden PRF key. The packed witness rows \(M\) and \(K\) introduced below are
only a representation of this abstract message on the witness domain. We
write the centering carry rows as \(J_0,J_1\).

The proof-system mechanics themselves are deferred to
Appendix~\ref{app:smallwood}. The purpose of the present section is narrower:
to fix the witness surfaces, the constraint families, and the replay basis for
issuance and showing.

\subsection{What SmallWood provides (formal interface)}
Fix a witness-support domain
\(\Omega=\{\omega_0,\dots,\omega_{s-1}\}\subset\Fq\) of size \(s\), a
disjoint tail set \(\Omega'\subset\Fq\setminus\Omega\), and a verifier opening
set \(E'\subset\Fq\setminus\Omega\) derived through Fiat--Shamir. A witness is
arranged as a matrix \(w\in\Fq^{n\times s}\), interpolated row-wise into
polynomials \(P_i(X)\in\Fq[X]\) such that
\[
  P_i(\omega_k)=w_{i,k}\qquad\text{for every }\omega_k\in\Omega,
\]
while the values on \(\Omega'\) are random HVZK tails. SmallWood commits to
the full row family under one PCS root and replays the active constraint
families from openings of those same committed rows at the verifier's points
\(E'\).

\paragraph{Parallel and aggregated families.}
SmallWood's PACS language distinguishes:
\begin{itemize}[leftmargin=1.25em]
  \item \emph{parallel families}, enforced pointwise on \(\Omega\), and
  \item \emph{aggregated families}, enforced through the identity
  \(\sum_{\omega\in\Omega}F'(\omega)=0\).
\end{itemize}
In our instantiation, parallel families capture slot-local algebra:
carrier membership, decode bridges, commitment, centering, packing, replay
residuals, and checkpointed PRF constraints. Aggregated families are used for
the public linear maps that bridge explicit-domain rows to replay-facing
transform aliases.

\paragraph{Replay basis.}
Linear relations are replayed directly from coefficient-basis openings. The
rational-hash relation is different: quotient-ring multiplication in
\(\Fq[X]/(X^N+1)\) is not preserved by naive scalar evaluation of coefficient
representatives at arbitrary points in \(E'\). For this reason, the prover also
commits transform-domain aliases for exactly those rows that participate in the
hash replay. The verifier checks the hash relation only in that transform
basis.

\subsection{Constraint compilation recipe}
Given a ring statement, our compilation follows five steps:
\begin{enumerate}[leftmargin=1.25em]
  \item \textbf{Choose the committed witness surface.}
  Commit explicit-domain witness rows for the low-alphabet carriers, the
  decoded aliases used in linear replay, and the minimal transform aliases
  required for quotient-ring replay.
  \item \textbf{Fix the packing on \(\Omega\).}
  Use coefficient packing: each coordinate of \(\Omega\) carries one packed
  coefficient slot.
  \item \textbf{Lift public rows into the replay basis.}
  Convert the public hash rows and public signature-map rows into
  theta-lifted explicit-domain polynomials, written
  \(B^\Theta\), \(T^\Theta\), and \(A^\Theta\).
  \item \textbf{Split the statement into parallel and aggregated families.}
  Slot-local algebra stays parallel; transform bridges are aggregated.
  \item \textbf{Provide formal coefficient polynomials.}
  In explicit-domain mode the verifier replays polynomial identities at
  \(E'\), so the prover supplies the formal coefficient descriptions of the
  active families.
\end{enumerate}

\subsection{Issuance (pre-sign) constraints}
\label{subsec:smallwood-issuance}
The pre-sign proof keeps the blind-signature target public and proves that it
is consistent with one hidden commitment opening and the issuer challenge. At
the ARC level, the statement binds together:
\begin{itemize}[leftmargin=1.25em]
  \item the hidden issuance witness
  \((\vecm,\vecr_0^{\Holder},\vecr_1^{\Holder},\bar\vecr)\),
  \item the public commitment \(\vecc\),
  \item the public issuer randomness
  \((\vecr_0^{\Issuer},\vecr_1^{\Issuer})\), and
  \item the public target \(T\) representing
  \(h_{\vecm,(\vecr_0,\vecr_1)}(B)\).
\end{itemize}

\paragraph{Committed witness rows.}
The pre-sign witness is committed in the following order:
\[
  C^M,\ C^{\mathrm{preRU}},\ C^{\mathrm{preR}},\ C^{\mathrm{ctr}},\ C^J,
\]
\[
  M,\ K,\ R_U^0,\ R_U^1,\ R,\ R_0,\ R_1,\ J_0,\ J_1,
\]
\[
  \widehat M,\ \widehat K,\ \widehat R_0,\ \widehat R_1.
\]
The first five rows are low-alphabet carriers. The next nine rows are decoded
aliases. The final four rows are the transform aliases used in the hash replay.

\paragraph{Carrier rows and decoded aliases.}
The message carrier \(C^M\) uses the sparse packed codec
\[
S_M^{\mathrm{packed}}
=
\{(a,0):a\in[-\BoundMsg,\BoundMsg]\}
\cup
\{(0,a):a\in[-\BoundMsg,\BoundMsg]\},
\]
with \((0,0)\) deduplicated and contiguous code set
\(\{0,\dots,4\BoundMsg\}\). The other carriers use the full pair codec
\[
  c=(x_0+B)+(2B+1)(x_1+B)
\]
with the appropriate bound for each family. In particular, \(C^J\) is a
bound-\(1\) pair carrier for \((J_0,J_1)\), and
\(C^{\mathrm{preR}}\) stores \((R,0)\).

The decoded aliases are determined by fixed public decode polynomials:
\[
  (M,K)=\mathrm{Dec}_M(C^M),
\]
\[
  (R_U^0,R_U^1)=\mathrm{Dec}_{\mathrm{preRU}}(C^{\mathrm{preRU}}),\qquad
  R=\mathrm{Dec}^{(1)}_{\mathrm{preR}}(C^{\mathrm{preR}}),
\]
\[
  (R_0,R_1)=\mathrm{Dec}_{\mathrm{ctr}}(C^{\mathrm{ctr}}),\qquad
  (J_0,J_1)=\mathrm{Dec}_J(C^J).
\]

\paragraph{Public inputs and lifted rows.}
The public issuance inputs are
\[
  \mACom,\ B,\ \vecc,\ \vecr_0^{\Issuer},\ \vecr_1^{\Issuer},\ T.
\]
The verifier also derives the lifts
\[
  B_0^\Theta,B_1^\Theta,B_2^\Theta,B_3^\Theta,\qquad T^\Theta,
\]
where \(T^\Theta\) is the theta-lift of the public coefficient-domain target.

\paragraph{Parallel issuance families.}
Let \(\Delta=2\BoundMsg+1\). The pre-sign proof enforces:
\begin{enumerate}[leftmargin=1.25em]
  \item \textbf{Carrier membership:}
  \[
    P_M(C^M)=0,\quad
    P_{\mathrm{preRU}}(C^{\mathrm{preRU}})=0,\quad
    P_{\mathrm{preR}}(C^{\mathrm{preR}})=0,\quad
    P_{\mathrm{ctr}}(C^{\mathrm{ctr}})=0,\quad
    P_J(C^J)=0.
  \]
  \item \textbf{Carrier-to-alias bridges:}
  \[
    M-\mathrm{Dec}_M^{(1)}(C^M)=0,\qquad
    K-\mathrm{Dec}_M^{(2)}(C^M)=0,
  \]
  \[
    R_U^0-\mathrm{Dec}_{\mathrm{preRU}}^{(1)}(C^{\mathrm{preRU}})=0,\qquad
    R_U^1-\mathrm{Dec}_{\mathrm{preRU}}^{(2)}(C^{\mathrm{preRU}})=0,
  \]
  \[
    R-\mathrm{Dec}_{\mathrm{preR}}^{(1)}(C^{\mathrm{preR}})=0,
  \]
  \[
    R_0-\mathrm{Dec}_{\mathrm{ctr}}^{(1)}(C^{\mathrm{ctr}})=0,\qquad
    R_1-\mathrm{Dec}_{\mathrm{ctr}}^{(2)}(C^{\mathrm{ctr}})=0,
  \]
  \[
    J_0-\mathrm{Dec}_J^{(1)}(C^J)=0,\qquad
    J_1-\mathrm{Dec}_J^{(2)}(C^J)=0.
  \]
  \item \textbf{Commitment:}
  \[
    \mACom\cdot[M\Vert K\Vert R_U^0\Vert R_U^1\Vert R]^\top=\vecc.
  \]
  \item \textbf{Centering:}
  \[
    R_0 = R_U^0 + \vecr_0^{\Issuer} - \Delta J_0,\qquad
    R_1 = R_U^1 + \vecr_1^{\Issuer} - \Delta J_1.
  \]
  \item \textbf{Packing:}
  \(M\) occupies the lower half of \(\Omega\), while \(K\) occupies the upper
  half, matching the packed representation of
  \(\vecm=\vecmu\Vert\veck\).
  \item \textbf{Transform-domain hash relation:}
  \[
    (B_3^\Theta-\widehat R_1)\,T^\Theta
    -\bigl(B_0^\Theta + B_1^\Theta(\widehat M+\widehat K)
    + B_2^\Theta\widehat R_0\bigr)=0.
  \]
\end{enumerate}

\paragraph{Aggregated issuance families.}
The issuance statement has four aggregated families:
\[
  M\rightarrow \widehat M,\qquad
  K\rightarrow \widehat K,\qquad
  R_0\rightarrow \widehat R_0,\qquad
  R_1\rightarrow \widehat R_1.
\]
No signature-side bridge appears in issuance, and there is no committed
\(\widehat T\) witness row because the target remains public.

\paragraph{Statement proved in issuance.}
The pre-sign proof asserts the existence of committed carriers, decoded aliases,
and non-sign transform aliases such that:
\begin{enumerate}[leftmargin=1.25em]
  \item each carrier lies in its carrier alphabet,
  \item each decoded alias is the decode image of its carrier,
  \item the decoded aliases satisfy commitment, centering, and packing,
  \item the transform aliases are the transform images of \(M,K,R_0,R_1\), and
  \item the public target \(T\), after theta-lift, satisfies the transform
  hash relation with those aliases.
\end{enumerate}

\subsection{Showing (post-sign) constraints}
\label{subsec:smallwood-showing}
The showing proof establishes knowledge of a valid signature on the hidden
message \(\vecm=\vecmu\Vert\veck\), hidden centered randomness
\((\vecr_0,\vecr_1)\), and a public tag satisfying
\(\prftag=\PRF(\veck,\nonce)\).

\paragraph{Public transform map.}
Let \(A^\Theta\) denote the theta-lifted public signature map and let
\(B_0^\Theta,B_1^\Theta,B_2^\Theta,B_3^\Theta\) denote the theta-lifted public
hash rows. The replay-facing transform aliases are related to their
explicit-domain source rows by one public linear map \(\mathcal F\), realized
through the bridge matrix
\[
  H_{t,k}=\mathrm{NTT}(e_k)_t,
\]
where \(e_k\) is the \(k\)-th coefficient basis vector.

\paragraph{Committed and virtual rows.}
The committed showing witness consists of:
\begin{enumerate}[leftmargin=1.25em]
  \item \textbf{Signature limb rows} \(L_{c,b,\ell}\),
  \item \textbf{Carrier rows} \(C^M\) and \(C^{\mathrm{ctr}}\),
  \item \textbf{Replay aliases}
  \(\widehat T_0,\widehat m,\widehat R_0,\widehat R_1\), where
  \(\widehat m=\mathcal F(M+K)\),
  \item \textbf{PRF companion rows.}
\end{enumerate}

The decoded non-sign rows remain virtual:
\[
  (M,K)=\mathrm{Dec}_M(C^M),\qquad
  (R_0,R_1)=\mathrm{Dec}_{\mathrm{ctr}}(C^{\mathrm{ctr}}).
\]
The packed signature source is reconstructed linearly from the limb surface:
\[
  \widetilde S_{c,b}=\sum_{\ell} r^\ell L_{c,b,\ell},
\]
and the intermediate signature hats are the corresponding transform images.

\paragraph{Parallel showing families.}
The showing proof enforces:
\begin{enumerate}[leftmargin=1.25em]
  \item \textbf{Signature digit membership:}
  the limb rows satisfy the signed low-alphabet digit constraints of the
  chosen shortness profile.
  \item \textbf{Replay residual:}
  \[
    (B_3^\Theta-\widehat R_1)\widehat T_0
    -\bigl(B_0^\Theta + B_1^\Theta\widehat m + B_2^\Theta\widehat R_0\bigr)=0.
  \]
  \item \textbf{Packing:}
  the decoded rows \(M,K\) satisfy the same half-split message
  representation as in issuance.
  \item \textbf{Carrier membership:}
  \[
    P_M(C^M)=0,\qquad P_{\mathrm{ctr}}(C^{\mathrm{ctr}})=0.
  \]
  \item \textbf{PRF key binding and nonlinear constraints:}
  the packed companion rows are bound to decoded \(K\)-coordinates, satisfy the
  grouped checkpoint equations, and bind the public tag after feed-forward and
  truncation.
\end{enumerate}

\paragraph{Aggregated showing families.}
The showing statement has three classes of aggregated families:
\begin{enumerate}[leftmargin=1.25em]
  \item \textbf{Direct limb-to-\(\widehat T_0\) bridge.}
  For each coordinate \(j\), one aggregated family proves that
  \(\widehat T_0(\omega_j)\) is the public linear image of the limb surface
  under the combined signature map and transform map.
  \item \textbf{Non-sign transform bridges.}
  The proof enforces
  \[
    (M+K)\rightarrow \widehat m,\qquad
    R_0\rightarrow \widehat R_0,\qquad
    R_1\rightarrow \widehat R_1.
  \]
  \item \textbf{PRF constancy for unpacked PRF rows.}
  This family belongs to the general SmallWood model but is empty for the
  packed companion representation used here.
\end{enumerate}

\paragraph{Statement proved in showing.}
The showing proof asserts the existence of committed limb rows, committed
carriers, committed replay aliases
\(\widehat T_0,\widehat m,\widehat R_0,\widehat R_1\), and PRF companion rows
such that:
\begin{enumerate}[leftmargin=1.25em]
  \item the cert-side signature is short via the limb decomposition,
  \item \(\widehat T_0\) is the public transform image of that short
  cert-side surface,
  \item \(\widehat m,\widehat R_0,\widehat R_1\) are the transform images of
  the decoded non-sign rows,
  \item the replay residual vanishes, and
  \item the PRF rows bind the public tag to the signed hidden key and the
  public nonce.
\end{enumerate}
Because \(\mathcal F\) is public and invertible, this is equivalent to the
coefficient-domain relation
\[
  B_3T - T\star R_1 - B_0 - B_1(M+K) - B_2R_0 = 0,
\]
together with the certification of \(T\) as the image of the short
cert-side signature witness.

\subsection{Poseidon PRF constraints with checkpoints}
\label{subsec:prf-checkpointed}
This subsection records the logical PRF constraints used in the showing
statement. The witness is carried through packed companion rows rather than one
row per logical lane, but the mathematical content is unchanged.

\paragraph{Logical PRF witness.}
Conceptually the PRF witness contains the hidden key lanes together with the
checkpointed nonlinear outputs. The packed companion representation stores these
logical values in packed rows and exposes the corresponding evaluations to the
verifier during replay.

\paragraph{Key binding.}
For each key lane, a public selector on \(\Omega\) binds the packed companion
value to the decoded packed-key row \(K\) at the designated coordinate. This
ties the hidden PRF key to the signed message component \(\veck\).

\paragraph{Checkpointed nonlinear constraints.}
The verifier recomputes the linear wiring between checkpoints and enforces the
grouped S-box equations at the committed checkpoints. These are parallel PACS
families. Under our parameters, the PRF degree remains below the signature
shortness degree and therefore does not dominate \(d_Q\).

\paragraph{Tag binding.}
The verifier reconstructs the final PRF state, applies the feed-forward with
the initial state, truncates to the public output lanes, and enforces equality
with the public tag.

\subsection{Where details are deferred}
Appendix~\ref{app:smallwood} records the SmallWood proof stack itself:
DECS/LVCS/PCS/PIOP, Fiat--Shamir, and the replay checks. The Gaussian trapdoor
sampler remains in Appendix~\ref{app:gaussian-sampler}. The PRF algorithm and
checkpoint schedule remain tied to Section~\ref{subsec:poseidon-prf} and
Appendix~\ref{app:prf}. The present section fixes only the statement layer.
